% Part: sets-functions-relations
% Chapter: functions
% Section: composition

\documentclass[../../../include/open-logic-section]{subfiles}

\begin{document}

\olfileid{sfr}{fun}{cmp}
\olsection{د تابعو ترکيب}

\begin{explain}
\oliflabeldef{sfr:fun:inv:sec}{په \olref[inv]{sec} کښې مو وليدل چې
د يوې \psOblique{!!a{bijection}}~$f$ معکوس~$f^{-1}$ پخپله يوه تابع ده۔ پر
تابعو يو بل عمل ترکيب دے: موږ}{موږ} دوه تابعې، $f$ او~$g$، سره
ترکيبولے او نوې تابع تعريفولے شو؛ يعنې لومړے~$f$ او بيا~$g$
تطبيق کړو۔ البته، دا يوازې هغه وخت ممکنه ده چې د قيمتونو سټونه
او د تعريف ساحې سره برابر راشي؛ يعنې د~$f$ د قيمتونو سټ بايد د~$g$
د تعريف د ساحې فرعي سټ وي۔ \oliflabeldef{sfr:rel:ops:sec}{پر تابعو
دا عمل، پر اړيکو د نسبي حاصل ضرب له هغه عمل سره ورته دے چې په
\olref[rel][ops]{sec} کښې راغلے دے۔}{}

انځوريز شکل د ترکيب د مفکورې په تشريح کښې مرسته کولے شي۔ په
\olref{fig:composition} کښې دوه تابعې $f \colon A \to B$ او
$g \colon B \to C$ او د هغو ترکيب~$(\comp{f}{g})$ ښيو۔ تابع
$(\comp{f}{g}) \colon A \to C$ د~$A$ هر !!{element} د~$C$ له يوه
\psOblique{!!a{element}} سره تړي۔ دا چې د $A$ يو !!a{element}
د~$C$ له کوم \psOblique{!!{element}} سره تړل کېږي، داسې ئې ټاکو:
د ورکړل شوي ننوت $x \in A$ دپاره، لومړے تابع $f$ پر~$x$ تطبيق
کړئ، چې يو $f(x) = y \in B$ به ورکړي؛ بيا تابع $g$ پر~$y$ تطبيق
کړئ، چې يو $g(f(x)) = g(y) = z \in C$ به ورکړي۔
\begin{figure}
  \olasset[2\olphotowidth]{assets/diagrams/composition.tikz}
  \caption{د دوو تابعو $f$ او~$g$ ترکيب $g \circ f$۔}
  \ollabel{fig:composition}
\end{figure}
\end{explain}

\begin{defn}[ترکيب]
فرض کړئ $f\colon A \to B$ او $g\colon B \to C$ تابعې دي۔
له~$g$ سره د $f$ \emph{ترکيب} دا دے: $\comp{f}{g} \colon A \to C$،
چې $(\comp{f}{g})(x) = g(f(x))$ وي۔
\end{defn}

\begin{ex}
تابعې $f(x) = x + 1$ او $g(x) = 2x$ په پام کښې ونيسئ۔ ځکه چې
$(\comp{f}{g})(x) = g(f(x))$، نو د هر ننوت~$x$ دپاره بايد لومړے
د هغه ځايناستے واخلئ، بيا پايله په دوو ضرب کړئ۔ نو د هغو ترکيب
په $(\comp{f}{g})(x) = 2(x+1)$ ورکړل شوے دے۔
\end{ex}

\begin{prob}
وښيئ چې که $f \colon A \to B$ او $g \colon B \to C$ دواړه
!!{injective} وي، نو $\comp{f}{g}\colon A \to C$ هم !!{injective} ده۔
\end{prob}

\begin{prob}
وښيئ چې که $f \colon A \to B$ او $g \colon B \to C$ دواړه
!!{surjective} وي، نو $\comp{f}{g}\colon A \to C$ هم !!{surjective} ده۔
\end{prob}

\begin{prob}
فرض کړئ $f \colon A \to B$ او $g \colon B \to C$۔ وښيئ چې د
$\comp{f}{g}$ ګراف $R_f \mid R_g$ دے۔
\end{prob}

\end{document}
